\documentclass[11pt]{scrartcl}
\usepackage[sexy]{evan}
\usepackage{answers}
\usepackage{hyperref}
\hypersetup{
    colorlinks=true,
    linkcolor=blue,
    urlcolor=teal
}

\begin{document}
\title{Mock Olympiad Solutions}
\author{Prakrit Gajurel}
\date{\today}
\renewcommand*\contentsname{Table of Contents}
\maketitle

\epigraph{As two clowns at the mercy of fate, let's fight this out
fair and square.}
{Natsuki Subaru}

Here's the set of solutions I've compiled for the SOMAES and MIN 
Physics and Math mock olympiad happening on August 16. These 
solutions are intended as a reference, and students may use different 
valid approaches. If something in a solution in this file is unclear, feel free 
to assume it and move on with the rest of the solution. In grading, however, any 
mathematically sound solution should receive full credit.



\tableofcontents

\newpage

\section{Official Solutions}

\subsection{Solution/s to Problem 1}

\begin{problem}
    Triangle $MIN$ is isosceles with $MI = MN$. Medians $\overline{IB}$ and
    $\overline{NA}$ are perpendicular to each other, and $IB = NA = 12$.
    Find, with proof, the area of $\triangle MIN$.
\end{problem}

\begin{remark}
    In what follows, we denote by $[X]$ the area contained within a polygon 
    $X$.
\end{remark}

\begin{figure}[htbp]
    \centering
    \includegraphics[width=0.5\textwidth]{amc10a problem.png}
    \caption{Diagram for the problem.}
\end{figure}

\subsubsection{Solution 1(areas)}

\begin{proof}
    We show that the answer is $96$. \\
    It's easy to observe that $[IABN] = 3[MBA]$, and also,
    \[
    [IABN] = \frac 12 (IB \cdot NA) = 72
    \]
    Giving $[MIN] = \dfrac 43 [IABN] = 96$.
\end{proof}

\subsubsection{Solution 2(centroid properties)}

\begin{proof}
    Let $G := IB \cap NA$ and let $C$ be the midpoint of $IN$.
    As $G$ is the centroid, it follows that
    $\dfrac{NG}{GA} = \dfrac 21$ giving $NG = 8$ and $GA = 4$,
    similarly we obtain $IG = 8 = 2GB$. Pythagoras on $\triangle AGI$
    gives $IA = AM = BM = NB = 4\sqrt5$, and Pythagoras on $\triangle IGN$
    gives $IN = 8\sqrt2 = 2CN$. Again, Pythagoras on $\triangle MNC$ gives 
    $MC = 12\sqrt2$, and we get our required answer as:
    \[
    [MIN] = \frac12 (MC \cdot IN) = \frac 12 (12\sqrt2 \cdot 8\sqrt2) = 96
    \]
    as desired.
\end{proof}

\subsubsection{Solution 3(cartesian bash)}

\begin{proof}
    We proceed in the Cartesian plane. Without loss of generality,
    let $I = (0,0)$ and $N = (n,0)$, giving $M = \left( \dfrac n2, h\right)$.
    Since $A$ and $B$ are the midpoints of $MI$ and $MN$, we obtain:
    \[
    A = \left( \frac n4, \frac h2 \right),
    \quad B = \left( \frac{3n}{4}, \frac h2 \right)
    \]
    Also, we obtain:
    \[
    \overrightarrow{IB} = \left\langle \frac{3n}{4}, \frac h2 \right\rangle,
    \quad \overrightarrow{NA} = \left\langle -\frac{3n}{4}, \frac h2 \right\rangle
    \]
    As $IB \perp NA$, we get that $h = \dfrac{3n}{2}$. Also using that $IB = 12$
    gives $n^2 = 128$ by the distance formula, giving:
    \[
    [MIN] = \frac 12 (nh) = \frac 12 \left( \frac{3n^2}{2} \right) = \boxed{96}.
    \]
\end{proof}

\newpage

\subsection{Solution/s to Problem 2}

\begin{problem}
    You are given that $x$ and $y$ are real numbers.

    \begin{enumerate}
        \item Given that $x+y$ and $x+y^2$ are rational numbers, is it always
        true that $x$ and $y$ are also rational numbers?

        \item Given that $x+y$, $x+y^2$, and $x+y^3$ are rational numbers,
        is it always true that $x$ and $y$ are also rational numbers?
    \end{enumerate}
\end{problem}

\subsubsection{Solution 1}

\begin{proof}
    Part $(1)$ is pretty straightforward. Obviously $\sqrt{p}$ for any prime
    $p$ is always irrational, and just considering something like 
    $x = -\dfrac{\sqrt p}{2}$ and $y = \dfrac{1+\sqrt{p}}{2}$ gives us that 
    the answer is no. \\ \\

    For part $(2)$, the answer is yes. Define the following:
    \[
    A = (x+y^2) - (x+y) = y(y-1), \qquad B = (x+y^3) - (x+y^2) = y^2(y-1)
    \]
    Obviously $A$ and $B$ are rational. When $A = 0$, we have $y(y-1)=0$ giving
    $y=0$ or $y=1$ giving $y$ rational, and so is $x = (x+y)-y$. When $A \neq 0$,
    $\dfrac{B}{A} = y$ is obviously rational, and of course, so is $x = (x+y)-y$. 
    Done.
\end{proof}

\newpage

\subsection{Solution/s to Problem 3}

\begin{problem}
    For every positive integer $n$, express, in terms of $n$, the number
    of $n$-digit positive integers satisfying both of the following conditions:

    \begin{itemize}
        \item No two consecutive digits are equal.
        \item The last digit is prime.
    \end{itemize}
\end{problem}

\subsubsection{Solution 1}

\begin{proof}
    Let $a_n$ denote the number of valid $n$-digit numbers whose last digit
    is prime. We claim that the answer is $a_n = \dfrac 25 (9^n-(-1)^n)$.
    We first begin by a counting argument. \\

    If we have the $n$-digit number, the last digit must be one of $2$, 
    $3$, $5$, or $7$. Let us consider a single one of these. Then, the 
    second last digit can be any number from $0$ to $9$ excluding the 
    prime digit we chose, giving us $9$ choices. Continuing this for the
    rest of the $n-1$ digits gives us an answer of $9^{n-1}$, but we had 
    $4$ choices for our primes so it is $4 \cdot 9^{n-1}$. We observe, however, 
    that the first digit could end up being $0$, giving us NOT the value of $a_n$
    but the value of $a_{n-1}$ and so to account for those cases, we add $a_{n-1}$ 
    to get our recursion $a_n + a_{n-1} = 4 \cdot 9^{n-1}$. \\

    Solving the recursion, we get:
    \begin{align*}
        a_n
        &= 4 \cdot 9^{n-1} - a_{n-1} \\
        &= 4 \cdot 9^{n-1} - 4 \cdot 9^{n-2} + a_{n-2} \\
        &= 4(9^{n-1} - 9^{n-2} + 9^{n-3} - \dots) \\
        &= 4 \left( \frac{9^n-(-1)^n}{10} \right) \\
        &= \frac25 (9^n-(-1)^n)
    \end{align*}
    We're done!
\end{proof}

\newpage

\subsection{Solution/s to Problem 4}

\begin{problem}
    Ram and Hari have created a program that,
    given a list of $n$ numbers
    $a_1, a_2, \dots, a_n$,
    calculates a list
    $b_1, b_2, \dots, b_n$
    such that, for each $1 \le k \le n$, the number $b_k$ is the number of
    times $a_k$ appears in the list. For example, the list
    \[
    1, \,2, \,3, \,1
    \]
    outputs
    \[
    2,\,1,\,1,\,2.
    \]

    Next, Ram asks Hari to run the program on a list of $2083$ numbers. He then
    asks him to apply the program to the resulting list, and so on, until a
    number greater than or equal to $k$ appears in the list. Find, with proof,
    the largest value of $k$ for which, whatever the initial list of $2083$
    positive integers is, it is always possible for Hari to do what Ram asked.
\end{problem}

\bigskip

\subsubsection{Solution 1}

\begin{proof}
    Define the $frequency$ of a number $x$ to be the number of times $x$ appears 
    in a list. \\

    We begin by proving the following claim:
    \begin{claim}
        The sequence of the frequencies for each position of the list is bounded 
        and non-decreasing.
        \begin{proof}
            We first show that the sequence is bounded. Since the list contains
            exactly $n=2083$ numbers, a number can appear at most $2083$ times.
            Also, a number must appear at least $1$ time in the list. Therefore,
            every number generated must be between $1$ and $2083$, giving us that
            our sequence is bounded. \\

            Consider in the first step, any number $x$ that appears $c$ times. 
            Then, in the list generated, the number $c$ will be present at least 
            $c$ times. The "at least" portion comes from the fact that multiple 
            distinct values in the original list could have the same frequency 
            $c$. Again, consider that number $c$ in the new list that appears $d$ 
            times. So, in the next list, all of the $c$'s turn into $d$'s, which 
            guarantees that the frequency of $d$ in the new list is $e$ where
            $e \ge d$. Therefore, the sequence of the frequencies of the elements 
            is non-decreasing.
        \end{proof}
    \end{claim}

    Since that sequence is bounded above and is non-decreasing, it must 
    eventually reach a constant state. Obviously, in that constant list, 
    any number $x$ will appear exactly $x$ times. We claim that the maximum 
    possible value of $k$ is $k=65$. \\ \\

    For our construction, just consider the following sequence:
    \[
    1, \,2, \,2, \,3, \,3, \,3, \, \dots \,, \, 65, \, 65
    \]
    excluding all $62$ of the $62$s. This construction obviously works as the 
    total number of items in the list is $\dfrac{65 \cdot 66}{2} - 62 = 2083$,
    so we're done for the construction. \\ \\

    We now show that $65$ is indeed the maximum. For the sake of contradiction,
    assume that we have a valid list for when it becomes constant such that all
    the numbers in the list are strictly less than $65$. Then, let the distinct numbers be 
    denoted by $n_1, \, n_2, \, \dots \,, \, n_{m}$ where $m \le 64$ assuming without loss 
    of generality that $n_1 < n_2 < \cdots < n_{m}$. This should give us that 
    $n_1 + n_2 + \cdots + n_{m} = 2083$, however,
    \[
    n_1 < n_2 < \cdots < n_{m} < 65 \implies
    n_1 + n_2 + \cdots + n_{m} \le 1 + 2 + \cdots + 64 < 2083,
    \]
    Contradiction! Thus, $65$ is the required answer.
\end{proof}

\end{document}
